% Zumdahl Chapter 2 free-response set -- 2 long (10) + 3 short (4) = 32 pts.
% New document type: a dedicated FRQ set per chapter, separate from the
% chapter test (which carries only 2 short FRQ of its own).
\documentclass{apchem-exam}

\apcunitword{Chapter}
\apcunit{2}
\apcunittitle{Atoms, Molecules, and Ions}
\apcdoctitle{Free-Response Set}
\apcsubtitle{Zumdahl \S2.2--\S2.8 \textbullet\ 5 questions \textbullet\ 32 points \textbullet\ 50 minutes}

\begin{document}
\apctitleblock

\begin{apcnote}[Scope --- read before assigning]
Only \textbf{CED topic 1.8} (Valence Electrons and Ionic Compounds,
Zumdahl \S2.6--\S2.8) is examinable material from this chapter. Question~1
is the one that maps directly onto the framework, and it is the one to
grade hardest.

Questions 2--5 rest on \S2.2--\S2.5 and \S2.8 --- the nuclear atom,
isotopes, the chemical laws, and nomenclature. The CED \emph{assumes}
this material rather than assessing it as standalone topics: the AP exam
will not ask you to recount the gold-foil experiment or to name a binary
compound for its own sake. They are here because every later chapter
depends on them, not because they will be tested in May. Treat them as
foundation work, not as exam rehearsal.
\end{apcnote}

\examsection{II}{Free Response}{5 questions \textbullet\ 32 points \textbullet\ 50 minutes}{\calculatorok}

% =====================================================================
% ---- FRQ 1 -- the CED-aligned one ------------------------------------------
\frq{long}{10}{Zumdahl \S2.6--\S2.8 \textbullet\ CED 1.8 \textbullet\ valence electrons, ionic compounds, Coulombic reasoning}

Magnesium and nitrogen react to form a single binary ionic compound.

\begin{parts}
  \item For magnesium and for nitrogen, state the group number, the number
        of valence electrons in the neutral atom, and the formula of the
        monatomic ion each element forms.
        \answerlines[3]{Mg: group 2, 2 valence electrons, forms
        \ce{Mg^2+}. \quad N: group 15, 5 valence electrons, forms
        \ce{N^3-}.}
  \item Write the formula of the compound they form, and show that your
        formula is electrically neutral. \answerlines[2]{\ce{Mg3N2};
        $3(+2) + 2(-3) = 0$.}
  \item Magnesium oxide and sodium chloride are both ionic solids.
        Predict which has the larger lattice energy, and justify your
        answer in terms of the Coulombic attraction between the ions.
        \workspace[3cm]
        \keyonly{\begin{apcanswer}\textbf{\ce{MgO}} has by far the larger
        lattice energy. Two Coulombic factors both push the same way.
        \textbf{Charge:} the attraction is proportional to the product of
        the ionic charges, which is $(+2)(-2) = 4$ for \ce{MgO} but
        $(+1)(-1) = 1$ for \ce{NaCl} --- a factor of four on its own.
        \textbf{Distance:} \ce{Mg^2+} and \ce{O^2-} are also the smaller
        pair, so the internuclear separation $d$ is shorter
        ($\approx$\SI{212}{\pico\metre} against
        $\approx$\SI{283}{\pico\metre}), and the attraction goes as
        $1/d$. Larger charge product over a shorter distance means a much
        stronger attraction and a much larger lattice
        energy.\end{apcanswer}}
  \item Magnesium oxide melts at \SI{2852}{\celsius}; sodium chloride
        melts at \SI{801}{\celsius}. Explain how this is consistent with
        your answer to (c). \answerlines[2]{Melting an ionic solid means
        overcoming the electrostatic attractions holding the lattice
        together. \ce{MgO} has the stronger attractions, so more energy
        is required and the melting point is much higher.}
  \item Solid \ce{MgO} does not conduct electricity, but molten \ce{MgO}
        does. Explain, and explain separately why an ionic crystal
        shatters when struck rather than bending. \workspace[3.2cm]
        \keyonly{\begin{apcanswer}\textbf{Conductivity.} Conduction
        requires charged particles that are free to move. In the solid the
        ions are locked in fixed lattice positions, so although charges
        are present none of them can migrate. Melting frees the ions to
        move through the liquid, so molten \ce{MgO} conducts.
        \textbf{Brittleness.} A blow displaces one plane of ions by one
        lattice position. Ions of \emph{like} charge are then brought face
        to face, the attraction along that plane is replaced by strong
        repulsion, and the crystal splits along it. A metal, whose
        delocalized electrons tolerate the same displacement, bends
        instead.\end{apcanswer}}
\end{parts}

\begin{rubric}
  \rpoint{2}{(a) 1 pt each element: group, valence count and ion all
    correct. The nitrogen ion must carry a $3-$ charge written as a
    superscript: \ce{N3-} is \emph{azide}, a different ion entirely, so
    insist on \ce{N^3-}}
  \rpoint{1}{(b) \ce{Mg3N2} with the neutrality shown --- the formula
    alone, unjustified, earns the point only if correct}
  \rpoint{2}{(c) 1 pt \ce{MgO} identified; 1 pt a Coulombic justification
    naming \textbf{charge} (and, for full credit at this level, distance).
    An answer resting on ``\ce{MgO} is more stable'', ``it obeys the
    octet rule'' or ``\ce{Mg} is more reactive'' earns \textbf{0} for this
    point --- no exceptions}
  \rpoint{2}{(d) 1 pt linking melting to breaking electrostatic
    attractions; 1 pt consistency with (c). Award the consistency point
    even if (c) was answered wrongly, provided the reasoning follows}
  \rpoint{3}{(e) 1 pt ions fixed in the solid; 1 pt ions mobile in the
    melt; 1 pt like charges brought adjacent producing repulsion. ``The
    bonds break'' earns 0 for the brittleness point}
\end{rubric}

% =====================================================================
% ---- FRQ 2 -----------------------------------------------------------------
\frq{long}{10}{Zumdahl \S2.3--\S2.5 \textbullet\ foundation \textbullet\ the nuclear atom and isotopes}

Chlorine occurs naturally as two isotopes, \ce{^{35}Cl} and
\ce{^{37}Cl}. The periodic table lists its atomic mass as
\SI{35.45}{\atomicmassunit}.

\begin{parts}
  \item In Rutherford's experiment a thin gold foil was bombarded with
        alpha particles. State the observation that was not expected, and
        state what it ruled out about the structure of the atom.
        \answerlines[3]{A very small fraction of the alpha particles were
        deflected through large angles, some almost straight back. This
        ruled out the ``plum pudding'' picture of charge and mass spread
        evenly through the atom: only a small, dense, positively charged
        centre could turn a fast alpha particle around.}
  \item Give the number of protons, neutrons and electrons in a neutral
        atom of \ce{^{37}Cl}, and in the ion \ce{^{37}Cl-}.
        \answerlines[3]{Neutral \ce{^{37}Cl}: 17 protons, 20 neutrons, 17
        electrons. \quad \ce{^{37}Cl-}: 17 protons, 20 neutrons, 18
        electrons --- the extra electron is what makes it an anion; the
        nucleus is untouched.}
  \item Explain why \ce{^{35}Cl} and \ce{^{37}Cl} have different masses
        but essentially identical chemical behaviour.
        \answerlines[3]{They differ only in neutron count, which changes
        the mass of the nucleus. Chemical behaviour is governed by the
        electrons --- their number and arrangement --- and both isotopes
        have 17 protons and therefore the same 17 electrons in the same
        configuration. Neutrons carry no charge and do not affect how the
        atom bonds.}
  \item Dalton's atomic theory states that all atoms of a given element
        are identical. Identify the part of that statement which the
        discovery of isotopes required chemists to modify, and state the
        corrected version. \answerlines[3]{The claim that atoms of an
        element are identical \emph{in mass} is wrong. The corrected
        statement is that all atoms of an element have the same number of
        protons (the same atomic number), but may differ in neutron
        number and therefore in mass. The rest of Dalton's theory
        survived.}
  \item Naturally occurring chlorine is \SI{75.8}{\percent} \ce{^{35}Cl}.
        Without calculating, state whether the average atomic mass lies
        closer to 35 or to 37, and justify your reasoning.
        \answerlines[2]{Closer to \textbf{35}. The average is weighted by
        abundance, and roughly three quarters of the atoms are the lighter
        isotope, so the mean is pulled toward 35 --- consistent with the
        tabulated \SI{35.45}{\atomicmassunit}.}
\end{parts}

\begin{rubric}
  \rpoint{2}{(a) 1 pt the large-angle/back-scattered deflection of a small
    fraction; 1 pt what it excludes (diffuse positive charge / plum
    pudding). ``It proved the nucleus exists'' alone earns 1, not 2}
  \rpoint{2}{(b) 1 pt the neutral atom fully correct; 1 pt the anion, which
    requires 18 electrons \emph{and} an unchanged nucleus. Changing the
    proton count loses the point}
  \rpoint{2}{(c) 1 pt different neutron number accounts for the mass
    difference; 1 pt chemistry is set by the electrons, which are the same}
  \rpoint{2}{(d) 1 pt identifying ``identical in mass'' as the failed
    clause; 1 pt the correction stated in terms of proton number}
  \rpoint{2}{(e) 1 pt ``closer to 35''; 1 pt an abundance-weighting
    justification. A numerical answer is acceptable but earns no extra
    credit --- the question asked for reasoning}
\end{rubric}

% =====================================================================
% ---- FRQ 3 -----------------------------------------------------------------
\frq{short}{4}{Zumdahl \S2.2 \textbullet\ foundation \textbullet\ the law of multiple proportions}

Two compounds contain only nitrogen and oxygen. Each sample below was
found to contain \SI{1.750}{\gram} of nitrogen.

\begin{center}
\begin{tabular}{lcc}
\hline
 & \textbf{Compound A} & \textbf{Compound B} \\
\hline
mass of nitrogen & \SI{1.750}{\gram} & \SI{1.750}{\gram} \\
mass of oxygen   & \SI{2.000}{\gram} & \SI{4.000}{\gram} \\
\hline
\end{tabular}
\end{center}

\begin{parts}
  \item Show that these data obey the law of multiple proportions.
        \workspace[2.4cm]
        \keyonly{\begin{apcanswer}Compare the oxygen combined with a
        \emph{fixed} mass of nitrogen. For A,
        $2.000/1.750 = \num{1.143}$~g O per g N; for B,
        $4.000/1.750 = \num{2.286}$~g O per g N. The ratio of these is
        $2.286/1.143 = \mathbf{2.00}$ --- a ratio of small whole numbers,
        which is exactly what the law requires. (Because the nitrogen mass
        is already the same in both samples, it is enough to note
        $4.000/2.000 = 2$.)\end{apcanswer}}
  \item Compound A is \ce{NO}. Deduce the formula of compound B, and state
        the reasoning. \answerlines[2]{\ce{NO2}. The same mass of nitrogen
        combines with twice as much oxygen in B, so B contains twice as
        many oxygen atoms per nitrogen atom.}
  \item State the postulate of Dalton's atomic theory that this result
        supports. \answerlines[2]{That compounds form when atoms combine
        in fixed, small whole-number ratios --- atoms are indivisible
        units, so a compound cannot contain a fractional atom.}
\end{parts}

\begin{rubric}
  \rpoint{2}{(a) 1 pt comparing oxygen masses at a fixed nitrogen mass;
    1 pt obtaining a small whole-number ratio of 2 (or 2:1). Dividing the
    two total sample masses instead earns 0 --- the comparison must hold
    nitrogen constant}
  \rpoint{1}{(b) \ce{NO2} with the doubling argument}
  \rpoint{1}{(c) any statement of combination in fixed small whole-number
    ratios}
\end{rubric}

% =====================================================================
% ---- FRQ 4 -----------------------------------------------------------------
\frq{short}{4}{Zumdahl \S2.7 \textbullet\ CED 1.8 \textbullet\ the periodic table and predicted ionic charge}

\begin{parts}
  \item Predict the charge of the monatomic ion formed by each of
        potassium, calcium, sulfur and bromine.
        \answerlines[2]{\ce{K+}, \ce{Ca^2+}, \ce{S^2-}, \ce{Br-}.}
  \item Sodium loses one electron readily, but removing a \emph{second}
        electron from sodium requires nearly ten times as much energy.
        Explain this in terms of the attraction experienced by the
        electron being removed. \workspace[3cm]
        \keyonly{\begin{apcanswer}The first electron removed is the single
        valence electron in the third shell. It is far from the nucleus
        and well shielded by the ten core electrons beneath it, so the net
        attraction holding it is small.

        Once it is gone, \ce{Na+} has a filled second shell. The next
        electron must come out of that \emph{core} shell --- much closer
        to the nucleus and shielded by only two electrons, so it
        experiences a far greater effective nuclear charge at a much
        shorter distance. Both Coulombic factors, larger $Z_{\text{eff}}$
        and smaller $r$, raise the energy required
        enormously.\end{apcanswer}}
\end{parts}

\begin{rubric}
  \rpoint{2}{(a) 1 pt for each correct pair; all four required for 2}
  \rpoint{2}{(b) 1 pt the second electron comes from a core shell closer
    to the nucleus; 1 pt greater effective nuclear charge / less
    shielding. An answer that says only ``\ce{Na+} has a full octet and is
    stable'' earns \textbf{0} --- the question asked for the attraction,
    not for a rule}
\end{rubric}

% =====================================================================
% ---- FRQ 5 -----------------------------------------------------------------
\frq{short}{4}{Zumdahl \S2.8 \textbullet\ foundation \textbullet\ naming simple compounds}

\begin{parts}
  \item Give the name of \ce{Fe2O3} and of \ce{N2O5}, and the formula of
        ammonium sulfate. \answerlines[2]{\ce{Fe2O3} is iron(III) oxide;
        \ce{N2O5} is dinitrogen pentoxide; ammonium sulfate is
        \ce{(NH4)2SO4}.}
  \item Iron compounds are named with a Roman numeral but sodium compounds
        are not. Explain why. \answerlines[3]{Sodium forms only one
        cation, \ce{Na+}, so its charge is never in doubt and no numeral
        is needed. Iron forms more than one stable cation --- \ce{Fe^2+}
        and \ce{Fe^3+} --- so the name would be ambiguous without stating
        which. The numeral gives the charge on the metal ion, not the
        number of atoms.}
\end{parts}

\begin{rubric}
  \rpoint{2}{(a) 2 pts all three correct; 1 pt for two. Note that
    \ce{N2O5} takes Greek prefixes because it is molecular, while
    \ce{Fe2O3} takes a Roman numeral because it is ionic --- a candidate
    who names \ce{N2O5} ``nitrogen(V) oxide'' earns no credit for it}
  \rpoint{2}{(b) 1 pt sodium has only one possible charge; 1 pt iron has
    more than one, so the numeral resolves the ambiguity. Saying the
    numeral counts atoms earns 0}
\end{rubric}

\examtotal

\end{document}
